%% Modelo mestre de currículo - layout dos PDFs de referência
%%
%% O agente deve adaptar o conteúdo à vaga sem alterar o preâmbulo,
%% as macros, o tamanho da página, as margens ou a ordem visual das seções.
%% Compilar documentos gerados com: lualatex main_<empresa>_<cargo>.tex

\documentclass[10pt,letterpaper]{article}
\usepackage[letterpaper,left=0.90in,right=0.45in,top=0.58in,bottom=0.60in]{geometry}
\usepackage[T1]{fontenc}
\usepackage[utf8]{inputenc}
\usepackage{mathptmx}
\usepackage[scaled=0.95]{helvet}
\usepackage{microtype}
\usepackage{tabularx}
\usepackage{enumitem}
\usepackage{xcolor}
\usepackage{hyperref}

\definecolor{rulegray}{RGB}{145,145,145}
\definecolor{contextgray}{RGB}{105,105,105}
\hypersetup{colorlinks=true,urlcolor=blue,linkcolor=blue,pdftitle={Candidate - CV},pdfpagemode=UseNone}
\pagestyle{empty}
\setlength{\parindent}{0pt}
\setlength{\parskip}{0pt}
\setlength{\emergencystretch}{2em}

% O cabeçalho/títulos usam Helvetica e o texto usa Times, como nos PDFs de referência.
\newcommand{\cvname}[1]{{\rmfamily\bfseries\fontsize{25}{28}\selectfont #1}\par\vspace{2pt}}
\newcommand{\contactline}[1]{{\sffamily\fontsize{9.5}{12}\selectfont #1}\par\vspace{12pt}}
\newcommand{\sectiontitle}[1]{\vspace{4pt}{\sffamily\bfseries\fontsize{15}{18}\selectfont #1}\par\vspace{1pt}{\color{rulegray}\rule{\linewidth}{0.45pt}}\par\vspace{5pt}}
\newcommand{\skillline}[2]{{\rmfamily\fontsize{10}{12}\selectfont\textbf{#1}: #2\par}}
\newcommand{\experienceentry}[3]{\begin{tabularx}{\linewidth}{@{}p{0.47\linewidth}@{}p{0.40\linewidth}@{}>{\raggedleft\arraybackslash}X@{}}{\rmfamily\bfseries #1} & {\rmfamily\itshape #2} & {\rmfamily #3}\\\end{tabularx}}
\newcommand{\educationentry}[3]{\begin{tabularx}{\linewidth}{@{}>{\bfseries}p{0.48\linewidth}@{}>{\raggedleft\arraybackslash}X@{}}#1 & #2\\\end{tabularx}\par{\rmfamily #3}\par}
\newenvironment{cvbullets}{\begin{itemize}[leftmargin=15pt,itemsep=1pt,topsep=0pt,parsep=0pt,partopsep=0pt]\rmfamily\fontsize{10}{12}\selectfont}{\end{itemize}}
\newcommand{\contextline}[1]{{\rmfamily\itshape\color{contextgray}\fontsize{9.5}{11.5}\selectfont Contexto: #1\par}}
\newcommand{\trainingline}[2]{{\rmfamily\fontsize{10}{12}\selectfont\textbf{#1}: #2\par}}

\begin{document}
\cvname{Candidate Name}
\contactline{\href{https://www.linkedin.com/in/PLACEHOLDER}{LinkedIn}\enspace | \enspace\href{https://github.com/PLACEHOLDER}{GitHub}\enspace | \enspace E-mail: YOUR.EMAIL@example.com\enspace | \enspace Telefone: +55 00 00000-0000}

\sectiontitle{Habilidades}
\skillline{Linguagens}{[LANGUAGE_1], [LANGUAGE_2], [LANGUAGE_3]}
\skillline{Backend e APIs}{[BACKEND_FRAMEWORK], REST APIs, [ORM_OR_DATA_ACCESS]}
\skillline{Frontend}{[FRONTEND_STACK]}
\skillline{Dados}{SQL, [DATABASE_1], [DATABASE_2]}
\skillline{Testes, DevOps e colaboração}{[TEST_STACK], Git, Docker, code review}
\skillline{Cloud}{[CLOUD_EXPERIENCE]}

\sectiontitle{Resumo}
{\rmfamily\fontsize{10}{12}\selectfont [PROFESSIONAL_SUMMARY_FOR_TARGET_ROLE]\par}

\sectiontitle{Experiências}
\experienceentry{[COMPANY_NAME_1]}{[JOB_TITLE_1]}{[MMM YYYY] - [MMM YYYY]}
\begin{cvbullets}
\item [IMPACT_BULLET_WITH_CONTEXT_AND_RESULT_1]
\item [IMPACT_BULLET_WITH_TECH_AND_SCOPE_2]
\item [IMPACT_BULLET_WITH_COLLABORATION_OR_DELIVERY_3]
\end{cvbullets}
\contextline{[SHORT_COMPANY_AND_DOMAIN_CONTEXT_1]}

\vspace{5pt}
\experienceentry{[COMPANY_NAME_2]}{[JOB_TITLE_2]}{[MMM YYYY] - [MMM YYYY]}
\begin{cvbullets}
\item [IMPACT_BULLET_1]
\item [IMPACT_BULLET_2]
\end{cvbullets}
\contextline{[SHORT_COMPANY_AND_DOMAIN_CONTEXT_2]}

\sectiontitle{Projeto em destaque}
{\rmfamily\fontsize{10}{12}\selectfont\textbf{[PROJECT_NAME] | [TECH_STACK] | }\href{https://github.com/PLACEHOLDER}{GitHub}\par}
[ONE_LINE_PROJECT_SUMMARY]
\begin{cvbullets}
\item [PROJECT_IMPLEMENTATION_DETAIL_1]
\item [PROJECT_ARCHITECTURE_OR_TEST_DETAIL_2]
\end{cvbullets}

\newpage
\sectiontitle{Formação}
\educationentry{[INSTITUTION_1]}{[MMM YYYY] - [MMM YYYY]}{[DEGREE_AND_FIELD_1]}
\educationentry{[INSTITUTION_2]}{[MMM YYYY] - [MMM YYYY]}{[DEGREE_AND_FIELD_2]}

\sectiontitle{Formação complementar}
\trainingline{[TOPIC_1]}{[COURSE_NAME_1] ([MMM/YYYY])}
\trainingline{[TOPIC_2]}{[COURSE_NAME_2] ([MMM/YYYY])}

\sectiontitle{Idiomas}
\trainingline{Português}{Nativo}
\trainingline{Inglês}{[ENGLISH_PROFICIENCY]}
\end{document}
